% Part:sets-functions-relations
% Chapter: sets
% Section: schroder-bernstein

\documentclass[../../../include/open-logic-section]{subfiles}

\begin{document}

\olfileid{sfr}{siz}{sb}

\olsection{আকারের ধারণা ও শ্র্যোডার--বার্নস্টাইন}

\begin{explain}
একটি স্বাভাবিক ধারণা এই: $A$,~$B$-এর আকারে বড় না হলে এবং $B$,~$A$-এর আকারে বড় না হলে $A$ ও $B$ সমসংখ্যক। সত্যি বলতে, এই ধারণাটি \emph{ভুল} হলে সমসংখ্যকতার যে ধারণা আমরা সংজ্ঞায়িত করেছি, সেটির সঙ্গে সেটের ``আকার'' তুলনার কোনো যোগ আছে—এ কথাই প্রায় সমর্থন করা যেত না!{} সৌভাগ্যক্রমে, স্বাভাবিক ধারণাটি সঠিক। শ্র্যোডার--বার্নস্টাইন উপপাদ্য এর সমর্থন দেয়।
\end{explain}

\begin{thm}[শ্র্যোডার--বার্নস্টাইন]
	\ollabel{thm:schroder-bernstein}
	যদি $\cardle{A}{B}$ ও $\cardle{B}{A}$ হয়,
	তবে $\cardeq{A}{B}$।
\end{thm}

\begin{explain}
অন্যভাবে বললে, $A$ থেকে~$B$-তে !!a{injection} এবং $B$ থেকে~$A$-তে !!a{injection} থাকলে $A$ থেকে~$B$-তে !!a{bijection} থাকে।

তবে এই ফলটি প্রমাণ করা সত্যিই বেশ \emph{কঠিন}। বস্তুত, কান্টর ফলটি বলেছিলেন, কিন্তু অন্যেরা তা প্রমাণ করেছিলেন।\footnote{ইতিহাস সম্পর্কে আরও জানতে, যেমন, \citet[pp.~165--6]{Potter2004} দেখো।}
\oliflabeldef{sfr:cardinals:card-sb:sec}{আমরা কেবল \olref[sfr][cardinals][card-sb]{sec}-এ শ্র্যোডার--বার্নস্টাইন \emph{প্রমাণ} করার অবস্থায় পৌঁছাব।}{}%
আপাতত ফলটি মেনে নিয়ে ব্যবহার করতে হবে।

সৌভাগ্যক্রমে, শ্র্যোডার--বার্নস্টাইন \emph{সঠিক}; তাই আমাদের সংজ্ঞায়িত সম্পর্কগুলি, অর্থাৎ $\cardeq{A}{B}$ ও $\cardle{A}{B}$, ``আকার''-এর সঙ্গে যুক্ত—এই চিন্তাকে এটি সমর্থন করে। তদুপরি, শ্র্যোডার--বার্নস্টাইন অত্যন্ত \emph{উপযোগী}। দুটি সমসংখ্যক সেটের মধ্যে !!a{bijection} ভাবা কঠিন হতে পারে। শ্র্যোডার--বার্নস্টাইন উপপাদ্য তুলনাটিকে দুটি ক্ষেত্রে ভাগ করতে দেয়: শুধু প্রথম সেট থেকে দ্বিতীয়টিতে এবং উল্টো দিকে !!a{injection} ভাবতে হয়।
\end{explain}
\end{document}
